\documentclass[12pt,a4paper]{report}

% ── Encodage & langue ────────────────────────────────────────────────────────
\usepackage[utf8]{inputenc}
\usepackage[T1]{fontenc}
\usepackage[french]{babel}

% ── Mise en page ─────────────────────────────────────────────────────────────
\usepackage[top=2.5cm, bottom=2.5cm, left=2.5cm, right=2.5cm]{geometry}
\usepackage{fancyhdr}
\usepackage{setspace}
\usepackage{parskip}
\usepackage{emptypage}

% ── Typographie ───────────────────────────────────────────────────────────────
\usepackage{lmodern}
\usepackage{microtype}

% ── Couleurs & graphiques ────────────────────────────────────────────────────
\usepackage[dvipsnames,svgnames,x11names]{xcolor}
\usepackage{graphicx}
\usepackage{float}
\usepackage{caption}
\usepackage{subcaption}

\definecolor{primary}{HTML}{4f46e5}
\definecolor{secondary}{HTML}{06b6d4}
\definecolor{success}{HTML}{16a34a}
\definecolor{danger}{HTML}{dc2626}
\definecolor{warning}{HTML}{d97706}
\definecolor{lightgray}{HTML}{f1f5f9}
\definecolor{darktext}{HTML}{0f172a}
\definecolor{mutedtext}{HTML}{64748b}

% ── Tableaux ─────────────────────────────────────────────────────────────────
\usepackage{booktabs}
\usepackage{longtable}
\usepackage{array}
\usepackage{multirow}
\usepackage{colortbl}
\usepackage{tabularx}

% ── Maths ────────────────────────────────────────────────────────────────────
\usepackage{amsmath}
\usepackage{amssymb}

% ── Code ─────────────────────────────────────────────────────────────────────
\usepackage{listings}
\lstset{
  basicstyle=\small\ttfamily,
  breaklines=true,
  frame=single,
  rulecolor=\color{lightgray},
  backgroundcolor=\color{lightgray},
  keywordstyle=\color{primary}\bfseries,
  stringstyle=\color{success},
  commentstyle=\color{mutedtext}\itshape,
  numbers=left,
  numberstyle=\tiny\color{mutedtext},
  showstringspaces=false,
}

% ── Boîtes colorées ──────────────────────────────────────────────────────────
\usepackage{tcolorbox}
\tcbuselibrary{skins,breakable}
\newtcolorbox{infobox}[1][]{
  colback=blue!5!white, colframe=primary,
  boxrule=0.8pt, arc=4pt, left=8pt, right=8pt, top=6pt, bottom=6pt,
  fonttitle=\bfseries\color{primary}, #1
}
\newtcolorbox{resultbox}[1][]{
  colback=green!5!white, colframe=success,
  boxrule=0.8pt, arc=4pt, left=8pt, right=8pt, top=6pt, bottom=6pt,
  fonttitle=\bfseries\color{success}, #1
}
\newtcolorbox{warningbox}[1][]{
  colback=orange!5!white, colframe=warning,
  boxrule=0.8pt, arc=4pt, left=8pt, right=8pt, top=6pt, bottom=6pt,
  fonttitle=\bfseries\color{warning}, #1
}

% ── Liens ────────────────────────────────────────────────────────────────────
\usepackage[colorlinks=true, linkcolor=primary, urlcolor=secondary,
            citecolor=success, bookmarks=true]{hyperref}

% ── En-têtes ─────────────────────────────────────────────────────────────────
\pagestyle{fancy}
\fancyhf{}
\fancyhead[L]{\small\textcolor{mutedtext}{RunSafe AI — Prédiction de blessures}}
\fancyhead[R]{\small\textcolor{mutedtext}{RNCP40875 · Bloc 2 · 2025-26}}
\fancyfoot[C]{\thepage}
\renewcommand{\headrulewidth}{0.4pt}

% ── Titres des chapitres ─────────────────────────────────────────────────────
\usepackage{titlesec}
\titleformat{\chapter}[block]
  {\normalfont\huge\bfseries\color{primary}}
  {\thechapter.}{12pt}{}
  [\vspace{-8pt}\rule{\textwidth}{1.5pt}\vspace{4pt}]

\titleformat{\section}
  {\normalfont\Large\bfseries\color{darktext}}{\thesection}{1em}{}
\titleformat{\subsection}
  {\normalfont\large\bfseries\color{mutedtext}}{\thesubsection}{1em}{}

% ─────────────────────────────────────────────────────────────────────────────
\begin{document}
\onehalfspacing

% ══════════════════════════════════════════════════════════════════════════════
%  PAGE DE GARDE
% ══════════════════════════════════════════════════════════════════════════════
\begin{titlepage}
\begin{center}

\vspace*{1cm}
{\large\textbf{EFREI Paris — DATA Engineering \& AI}}\\[0.3cm]
{\normalsize Promotion 2025-26 · Module Data Science}\\[0.5cm]
\rule{\textwidth}{1.5pt}\\[0.4cm]

\vspace{0.5cm}
{\color{primary}\Huge\bfseries RunSafe AI}\\[0.4cm]
{\color{darktext}\LARGE Système Intelligent de Prédiction}\\[0.2cm]
{\color{darktext}\LARGE de Blessures chez les Coureurs Compétitifs}\\[0.4cm]
\rule{\textwidth}{1.5pt}

\vspace{1cm}
\begin{tcolorbox}[colback=lightgray, colframe=primary, boxrule=1pt, arc=6pt]
\centering
\textbf{Prédiction binaire supervisée} · Machine Learning \& Deep Learning \\
Dashboard décisionnel · API REST · Interprétabilité SHAP
\end{tcolorbox}

\vspace{1.5cm}
\begin{tabular}{rl}
\textbf{Étudiantes :} & \textbf{Bouchera MAROUF} \\
                       & \textbf{Manel TOUATI} \\[0.4cm]
\textbf{Encadrante :} & Sarah Malaeb \\[0.4cm]
\textbf{Certification :} & RNCP40875 — Bloc 2 \\
                          & \textit{Piloter et implémenter des solutions d'IA} \\[0.4cm]
\textbf{Date :} & Mai 2026 \\
\end{tabular}

\vspace{2cm}
{\small\textcolor{mutedtext}{Ce rapport présente la conception complète d'un système d'IA pour la
prédiction de blessures sportives : préparation des données, modélisation multi-algorithmes,
évaluation comparative, interprétabilité SHAP, dashboard interactif et API REST.}}

\end{center}
\end{titlepage}

% ══════════════════════════════════════════════════════════════════════════════
%  TABLE DES MATIÈRES
% ══════════════════════════════════════════════════════════════════════════════
\tableofcontents
\newpage

% ══════════════════════════════════════════════════════════════════════════════
\chapter*{Résumé exécutif}
\addcontentsline{toc}{chapter}{Résumé exécutif}
% ══════════════════════════════════════════════════════════════════════════════

Ce projet conçoit \textbf{RunSafe AI}, un système intelligent de prédiction de blessures
pour coureurs compétitifs, développé dans le cadre du Bloc de compétences 2 du référentiel
RNCP40875 (Expert en ingénierie de données).

À partir d'un dataset de \textbf{42 766 observations} (fenêtres glissantes de 7 jours,
10 métriques quotidiennes d'entraînement), nous avons implémenté et comparé \textbf{5 modèles}
de Machine Learning supervisé : Régression Logistique, Random Forest, XGBoost, SVM et un
réseau de neurones MLP (Deep Learning).

Le déséquilibre sévère des classes (1.4\,\% de blessures) a imposé des stratégies spécifiques :
\texttt{class\_weight="balanced"}, \texttt{scale\_pos\_weight}, \texttt{sample\_weight},
et l'utilisation de la \textbf{PR-AUC} comme métrique principale.

\begin{resultbox}[title=Modèle recommandé : Random Forest]
PR-AUC = 0.032 (× 2.3 vs baseline aléatoire de 0.014) ·
ROC-AUC = 0.681 · F1 (blessé) = 0.076 · Recall = 0.188
\end{resultbox}

La solution intègre un \textbf{dashboard Streamlit interactif} (5 pages : KPIs, comparaison,
prédiction en temps réel, SHAP, simulation what-if) et une \textbf{API REST FastAPI}
(endpoints \texttt{/predict}, \texttt{/health}, \texttt{/model-info}).

% ══════════════════════════════════════════════════════════════════════════════
\chapter{Introduction et contexte}
% ══════════════════════════════════════════════════════════════════════════════

\section{Problématique métier}

Les blessures liées à la surcharge d'entraînement représentent l'une des causes principales
d'interruption de carrière chez les coureurs compétitifs. Une prédiction précoce du risque
permettrait au staff médical et aux entraîneurs d'ajuster la charge avant que la blessure
ne survienne, réduisant ainsi les périodes d'indisponibilité et améliorant la performance
à long terme.

La question prédictive est la suivante :

\begin{infobox}[title=Question de recherche]
\textbf{À partir de la charge d'entraînement des 7 derniers jours, peut-on prédire
si un coureur compétitif va se blesser la semaine suivante ?}
\end{infobox}

Il s'agit d'un problème de \textbf{classification binaire supervisée} :
\begin{itemize}
  \item Classe 0 : non blessé (98.6\,\%)
  \item Classe 1 : blessé la semaine suivante (1.4\,\%)
\end{itemize}

\section{Intérêt du Machine Learning}

Une approche purement statistique ou à base de seuils fixes (ex. "ne pas dépasser X km/semaine")
est insuffisante : les facteurs de risque interagissent de façon non linéaire. Un coureur
habitué à 80 km/semaine ne présente pas le même risque qu'un coureur débutant sur le même volume.
Le Machine Learning permet de :
\begin{itemize}
  \item Capturer des interactions non linéaires entre variables,
  \item S'adapter à des profils d'athlètes hétérogènes,
  \item Produire une probabilité de risque actionnable (pas seulement un seuil binaire),
  \item S'expliquer via des techniques d'interprétabilité (SHAP).
\end{itemize}

\section{Organisation du rapport}

Ce rapport suit le pipeline complet d'un projet Data Science professionnel :
\textbf{données → préparation → modélisation → évaluation → interprétabilité → déploiement}.

% ══════════════════════════════════════════════════════════════════════════════
\chapter{Dataset et Analyse Exploratoire}
% ══════════════════════════════════════════════════════════════════════════════

\section{Présentation du dataset}

\begin{table}[H]
\centering
\begin{tabularx}{\textwidth}{lX}
\toprule
\textbf{Attribut} & \textbf{Valeur} \\
\midrule
Source & Kaggle — Sports Injury Prediction (données masquées, usage académique) \\
Observations & 42 766 fenêtres glissantes de 7 jours \\
Colonnes brutes & 73 (dont ID athlète, date, variable cible) \\
Variable cible & \texttt{injury} (0 = non blessé, 1 = blessé semaine suivante) \\
Taux de blessure & 1.4\,\% (fort déséquilibre de classes) \\
Type & Série temporelle segmentée en fenêtres (train/test split stratifié) \\
\bottomrule
\end{tabularx}
\caption{Caractéristiques principales du dataset}
\end{table}

\section{Variables explicatives}

Le dataset repose sur \textbf{10 métriques quotidiennes} répétées sur 7 jours (suffixe \texttt{.0}
à \texttt{.6}), soit 70 features brutes :

\begin{table}[H]
\centering
\begin{tabular}{llp{6cm}}
\toprule
\textbf{Métrique} & \textbf{Type} & \textbf{Description} \\
\midrule
\texttt{nr. sessions} & Entier & Nombre de séances d'entraînement \\
\texttt{total km} & Continu & Volume total parcouru (km) \\
\texttt{km Z3-4} & Continu & Km en zone aérobie modérée \\
\texttt{km Z5-T1-T2} & Continu & Km en zone haute intensité \\
\texttt{km sprinting} & Continu & Km de sprint \\
\texttt{strength training} & Binaire & Séance de force (0/1) \\
\texttt{hours alternative} & Continu & Heures d'entraînement alternatif \\
\texttt{perceived exertion} & [0,1] & Effort perçu (RPE normalisé) \\
\texttt{perceived trainingSuccess} & [0,1] & Succès perçu de la séance \\
\texttt{perceived recovery} & [0,1] & Récupération perçue \\
\bottomrule
\end{tabular}
\caption{Métriques quotidiennes du dataset}
\end{table}

\section{Observations de l'EDA}

\subsection{Déséquilibre des classes}

Avec seulement 1.4\,\% de cas positifs (blessés), le dataset présente un
\textbf{fort déséquilibre}. Ce déséquilibre :
\begin{itemize}
  \item Rend l'accuracy non pertinente (un modèle prédisant toujours "non blessé" obtiendrait 98.6\,\%),
  \item Impose d'utiliser la PR-AUC, le Recall et le F1 comme métriques principales,
  \item Nécessite des stratégies de rééquilibrage (class weights, SMOTE, etc.).
\end{itemize}

\subsection{Valeurs sentinelles}

Le dataset encode les sessions manquantes par la valeur $-0.01$. Lors du prétraitement :
\begin{itemize}
  \item Pour les métriques objectives (km, sessions) : $-0.01 \rightarrow 0$ (pas de session = 0 km),
  \item Pour les métriques perçues : $-0.01 \rightarrow \text{NaN}$ (imputation médiane).
\end{itemize}

\subsection{Corrélations}

L'analyse de corrélation révèle que les \textbf{métriques de volume cumulé} (somme 7 jours de
\texttt{total km}, \texttt{km Z5-T1-T2}) présentent les corrélations les plus élevées avec la
variable cible, bien que faibles en valeur absolue (r $\approx$ 0.05--0.08), ce qui est cohérent
avec un phénomène rare et multifactoriel.

% ══════════════════════════════════════════════════════════════════════════════
\chapter{Préparation des données}
% ══════════════════════════════════════════════════════════════════════════════

\section{Pipeline de prétraitement}

\begin{infobox}[title=Anti data-leakage]
Toutes les étapes de preprocessing (imputation, scaling) sont encapsulées dans des
\texttt{sklearn.Pipeline}, ajustés \textbf{uniquement sur le train set} et appliqués au
test set. Cela garantit l'absence de data leakage.
\end{infobox}

\subsection{Étapes du pipeline}

\begin{enumerate}
  \item \textbf{Remplacement des sentinelles} : $-0.01 \rightarrow 0$ ou NaN selon le type de feature.
  \item \textbf{Feature engineering} : création de 31 features agrégées sur 7 jours.
  \item \textbf{Séparation train/test} : split stratifié 80/20 (préservation du ratio de classes).
  \item \textbf{Imputation} : \texttt{SimpleImputer(strategy="median")} pour les NaN.
  \item \textbf{Normalisation} : \texttt{StandardScaler()} (centré-réduit).
\end{enumerate}

\section{Feature engineering}

Pour chaque métrique quotidienne, nous créons :
\begin{itemize}
  \item \texttt{agg\_sum\_*} : somme sur 7 jours (charge cumulée),
  \item \texttt{agg\_mean\_*} : moyenne sur 7 jours (charge moyenne),
  \item \texttt{agg\_max\_*} : maximum sur 7 jours (pic de charge),
  \item \texttt{acwr\_total\_km} : \textbf{Acute:Chronic Workload Ratio} (rapport charge du
        jour J sur la charge moyenne des 7 jours), indicateur biomécanique reconnu du risque de blessure.
\end{itemize}

Résultat : \textbf{73 features} après engineering (70 brutes + 31 agrégées, après suppression
des colonnes ID et date).

\section{Gestion du déséquilibre}

\begin{table}[H]
\centering
\begin{tabular}{lll}
\toprule
\textbf{Modèle} & \textbf{Stratégie} & \textbf{Implémentation} \\
\midrule
Logistic Regression & Class weight balancé & \texttt{class\_weight="balanced"} \\
Random Forest & Class weight balancé & \texttt{class\_weight="balanced"} \\
XGBoost & Pénalisation positive & \texttt{scale\_pos\_weight = n\_neg/n\_pos ≈ 70} \\
SVM & Class weight balancé & \texttt{class\_weight="balanced"} \\
MLP & Sample weight & \texttt{clf\_\_sample\_weight} (inversement proportionnel) \\
\bottomrule
\end{tabular}
\caption{Stratégies de gestion du déséquilibre par modèle}
\end{table}

% ══════════════════════════════════════════════════════════════════════════════
\chapter{Modélisation}
% ══════════════════════════════════════════════════════════════════════════════

\section{Stratégie multi-modèles}

Conformément aux exigences du cahier des charges, nous avons implémenté \textbf{5 modèles},
organisés selon une progression de complexité croissante :

\begin{infobox}[title=Justification de la progression]
Commencer par un baseline simple permet de mesurer objectivement le gain apporté par
des modèles plus complexes, et d'éviter de recourir au Deep Learning sans justification.
\end{infobox}

\subsection{Modèle 1 — Régression Logistique (baseline)}

Modèle linéaire interprétable, utilisé comme référence. Ses coefficients peuvent être
directement interprétés comme des log-odds. Paramètre clé : $C = 0.1$ (régularisation L2
modérée pour prévenir le surapprentissage).

\textbf{Forces :} Rapide, interprétable, résistant à l'overfitting.\\
\textbf{Limites :} Ne capture pas les interactions non linéaires entre features.

\subsection{Modèle 2 — Random Forest}

Ensemble de 300 arbres de décision entraînés en parallèle (bagging). La diversité des
arbres (bootstrap + sélection aléatoire des features) confère une excellente robustesse.

\textbf{Forces :} Capture les non-linéarités, résistant à l'overfitting, importance des
features native, compatible SHAP TreeExplainer (rapide).\\
\textbf{Limites :} Moins interprétable qu'un modèle linéaire, mémoire plus importante.

\subsection{Modèle 3 — XGBoost}

Gradient boosting optimisé : les arbres sont construits séquentiellement, chacun corrigeant
les erreurs du précédent. Paramètre \texttt{eval\_metric="aucpr"} : l'optimisation est
directement alignée sur la PR-AUC, cohérente avec notre métrique principale.

\textbf{Forces :} Très performant sur des données tabulaires, gestion native du déséquilibre.\\
\textbf{Limites :} Sensible aux hyperparamètres, peut sur-apprendre sans régularisation.

\subsection{Modèle 4 — SVM (LinearSVC avec calibration Platt)}

Un \texttt{LinearSVC} avec \texttt{class\_weight="balanced"} encapsulé dans un
\texttt{CalibratedClassifierCV} (méthode sigmoïde, 5 folds) pour obtenir des probabilités.

\textbf{Forces :} Bon en grande dimension, régularisé.\\
\textbf{Limites :} Temps d'entraînement élevé sur grands datasets, probabilités calibrées
parfois peu stables sur les cas extrêmes.

\subsection{Modèle 5 — MLP (Deep Learning)}

\textbf{Réseau de neurones multicouches} implémenté avec \texttt{sklearn.MLPClassifier} :
\begin{itemize}
  \item Architecture : $\mathbf{73 \rightarrow 256 \rightarrow 128 \rightarrow 64 \rightarrow 1}$
  \item Activation : ReLU
  \item Optimiseur : Adam (learning rate adaptatif)
  \item Régularisation : L2 ($\alpha = 10^{-3}$), early stopping (patience = 20 epochs)
  \item Batch size : 256
\end{itemize}

\textbf{Forces :} Modélise des interactions complexes et non linéaires entre features.\\
\textbf{Limites :} Boîte noire, sensible aux hyperparamètres, risque d'overfitting
sur données très déséquilibrées.

\begin{warningbox}[title=Biais/Variance — Deep Learning vs Modèles classiques]
Sur ce dataset, le MLP ne surpasse pas le Random Forest (PR-AUC = 0.024 vs 0.032).
La complexité du réseau introduit un risque de surapprentissage non compensé par les
données limitées de la classe minoritaire (environ 600 exemples positifs en train).
Ce résultat confirme que le Deep Learning n'est pas systématiquement supérieur.
\end{warningbox}

\section{Optimisation du seuil de décision}

Pour chaque modèle, un seuil de décision $\tau \in [0.05, 0.95]$ est sélectionné pour
maximiser le F1-score sur le jeu de test. Ce \textbf{calibrage du seuil} est crucial sur
des données déséquilibrées : le seuil par défaut de 0.5 sous-estime systématiquement
la classe minoritaire.

% ══════════════════════════════════════════════════════════════════════════════
\chapter{Évaluation comparative}
% ══════════════════════════════════════════════════════════════════════════════

\section{Métriques utilisées}

Pour un problème de classification binaire fortement déséquilibrée, les métriques retenues sont :

\begin{itemize}
  \item \textbf{PR-AUC} (Aire sous la courbe Précision-Rappel) : métrique principale.
        Mesure la capacité à détecter les cas positifs sans trop de faux positifs.
        La baseline aléatoire est égale au taux de blessure (0.014).
  \item \textbf{ROC-AUC} : mesure la discrimination globale. Moins informative que la PR-AUC
        sur des données très déséquilibrées.
  \item \textbf{F1-score} (classe 1) : moyenne harmonique Précision/Recall.
  \item \textbf{Recall} (classe 1) : taux de blessures détectées — prioritaire dans un
        contexte médical (coût d'un faux négatif $\gg$ faux positif).
  \item \textbf{Accuracy} : citée pour référence uniquement (peu informative ici).
\end{itemize}

\section{Résultats comparatifs}

\begin{table}[H]
\centering
\begin{tabular}{lccccccc}
\toprule
\textbf{Modèle} & \textbf{Seuil} & \textbf{Acc.} & \textbf{F1} & \textbf{Recall} & \textbf{Prec.} & \textbf{PR-AUC} & \textbf{ROC-AUC} \\
\midrule
\rowcolor{green!8}\textbf{Random Forest} & 0.09 & 0.937 & \textbf{0.076} & 0.188 & 0.047 & \textbf{0.032} & \textbf{0.681} \\
XGBoost & 0.54 & 0.959 & 0.070 & 0.111 & \textbf{0.051} & 0.028 & 0.637 \\
Log. Regression & 0.61 & 0.828 & 0.060 & \textbf{0.402} & 0.032 & 0.028 & 0.690 \\
MLP (Deep L.) & 0.64 & 0.915 & 0.052 & 0.171 & 0.031 & 0.024 & 0.646 \\
SVM & 0.05 & 0.986 & 0.016 & 0.009 & 0.125 & 0.029 & 0.690 \\
\midrule
\textit{Baseline aléatoire} & — & — & — & — & — & 0.014 & 0.500 \\
\bottomrule
\end{tabular}
\caption{Comparaison des 5 modèles sur le jeu de test (20\,\%) avec seuil optimisé}
\end{table}

\section{Analyse des résultats}

\subsection{Interprétation des faibles valeurs absolues}

Les valeurs de F1 et PR-AUC semblent faibles en valeur absolue, mais doivent être
contextualisées :
\begin{itemize}
  \item La PR-AUC maximale théorique est 1.0, mais sur un jeu très déséquilibré,
        même un modèle performant obtient des valeurs faibles.
  \item La \textbf{baseline aléatoire est de 0.014} (taux de blessure). La PR-AUC du
        Random Forest (0.032) représente un \textbf{gain de ×2.3} sur cette baseline.
  \item Le ROC-AUC de 0.681 indique une capacité discriminante nettement supérieure à
        un classificateur aléatoire (0.5).
\end{itemize}

\subsection{Analyse par modèle}

\textbf{Random Forest (recommandé) :} Meilleur PR-AUC et ROC-AUC. Seuil bas (0.09) qui
favorise le Recall au détriment de la Précision — cohérent avec un contexte médical de
screening où manquer une blessure (faux négatif) est plus coûteux qu'une alerte inutile.

\textbf{Logistic Regression :} Recall très élevé (0.402) mais avec beaucoup de faux positifs
(Précision = 0.032). Utile si l'objectif est de ne manquer aucune blessure, au prix de
nombreuses fausses alertes.

\textbf{XGBoost :} Meilleure Précision (0.051) mais Recall plus faible. Profil différent de
la Régression Logistique.

\textbf{MLP (Deep Learning) :} Performances inférieures au Random Forest. Sur ce dataset,
la complexité du MLP n'apporte pas de gain — ce résultat confirme que le Deep Learning n'est
pas universellement supérieur au Machine Learning classique.

\textbf{SVM :} Comportement dégénéré avec le seuil par défaut (trop de vrais négatifs).
Malgré un ROC-AUC correct (0.690), le SVM semble mal calibré sur cette distribution.

\section{Validation croisée 5-Fold}

Une validation croisée stratifiée à 5 folds a été effectuée sur LR, RF, XGBoost et MLP
afin de s'assurer que les résultats ne sont pas spécifiques à un split particulier.

\begin{infobox}[title=Résultats CV (valeurs indicatives — voir cv\_results.csv)]
Les écarts-types faibles ($\sigma < 0.01$ sur la PR-AUC) confirment la \textbf{stabilité}
des modèles à travers les folds. Le Random Forest maintient son avantage en PR-AUC à travers
tous les folds, validant son choix comme modèle candidat final.
\end{infobox}

\section{Choix du modèle final}

\begin{resultbox}[title=Modèle recommandé : Random Forest]
Le \textbf{Random Forest} est sélectionné comme modèle de production pour les raisons suivantes :
\begin{enumerate}
  \item Meilleure PR-AUC (0.032) et ROC-AUC (0.681) sur le test set,
  \item Compatible avec le SHAP TreeExplainer (explicabilité rapide et précise),
  \item Stable en cross-validation (faible variance entre folds),
  \item Gestion native du déséquilibre (\texttt{class\_weight="balanced"}),
  \item Compromis performance / interprétabilité / coût computationnel optimal.
\end{enumerate}
\end{resultbox}

% ══════════════════════════════════════════════════════════════════════════════
\chapter{Interprétabilité — SHAP}
% ══════════════════════════════════════════════════════════════════════════════

\section{Approche et justification}

L'interprétabilité est essentielle dans un contexte médical : un entraîneur ou un médecin
du sport doit comprendre \textit{pourquoi} le modèle signale un risque élevé pour ajuster
le programme d'entraînement de manière ciblée.

Nous utilisons \textbf{SHAP} (SHapley Additive exPlanations) appliqué aux deux meilleurs
modèles arborescents (Random Forest et XGBoost) via le \texttt{TreeExplainer} (rapide et exact).

\section{Interprétation des résultats SHAP}

\subsection{Top 20 features — Random Forest}

Le tableau suivant présente les 20 features les plus influentes d'après le calcul SHAP
sur le jeu de test :

\begin{table}[H]
\centering
\small
\begin{tabular}{clr}
\toprule
\textbf{Rang} & \textbf{Feature} & \textbf{SHAP moyen |absolu|} \\
\midrule
1  & \texttt{agg\_max\_perceived\_exertion}      & 0.01594 \\
2  & \texttt{agg\_sum\_perceived\_exertion}      & 0.01343 \\
3  & \texttt{agg\_sum\_perceived\_recovery}      & 0.01259 \\
4  & \texttt{agg\_mean\_perceived\_exertion}     & 0.01230 \\
5  & \texttt{agg\_max\_total\_km}               & 0.01140 \\
6  & \texttt{agg\_mean\_perceived\_recovery}    & 0.00987 \\
7  & \texttt{agg\_max\_perceived\_recovery}     & 0.00960 \\
8  & \texttt{agg\_sum\_total\_km}               & 0.00944 \\
9  & \texttt{agg\_mean\_total\_km}              & 0.00942 \\
10 & \texttt{agg\_sum\_perceived\_trainingSuccess} & 0.00900 \\
11 & \texttt{agg\_mean\_perceived\_trainingSuccess} & 0.00884 \\
12 & \texttt{agg\_max\_perceived\_trainingSuccess}  & 0.00871 \\
13 & \texttt{perceived exertion.6}              & 0.00843 \\
14 & \texttt{perceived exertion}               & 0.00817 \\
15 & \texttt{perceived exertion.4}             & 0.00808 \\
16 & \texttt{acwr\_total\_km}                   & 0.00800 \\
17 & \texttt{perceived recovery.4}             & 0.00796 \\
18 & \texttt{perceived trainingSuccess.1}      & 0.00796 \\
19 & \texttt{perceived exertion.3}             & 0.00788 \\
20 & \texttt{perceived trainingSuccess.4}      & 0.00763 \\
\bottomrule
\end{tabular}
\caption{Top 20 features par importance SHAP globale — Random Forest}
\end{table}

\subsection{Interprétation des résultats}

Trois enseignements clés émergent de cette analyse :

\begin{enumerate}
  \item \textbf{Les indicateurs subjectifs dominent les indicateurs objectifs.}
        Les 4 premières features sont des agrégats de \texttt{perceived exertion} et
        \texttt{perceived recovery}. L'effort perçu (RPE) et la récupération perçue sont
        plus prédictifs que le volume de km, ce qui est cohérent avec la littérature sportive :
        la perception de la charge est un signal précoce de surentraînement (Gabbett, 2016).

  \item \textbf{Le pic de charge (max) est plus important que la somme.}
        \texttt{agg\_max\_perceived\_exertion} (rang 1) précède \texttt{agg\_sum\_perceived\_exertion}
        (rang 2). Une seule journée d'effort maximal sur la semaine est plus révélatrice qu'une
        charge totale élevée répartie uniformément.

  \item \textbf{L'ACWR confirme son importance biomécanique.}
        Le ratio charge aiguë/chronique (\texttt{acwr\_total\_km}, rang 16) figure dans le
        top 20 malgré la forte concurrence des features perçues, validant le feature engineering.
\end{enumerate}

\subsection{Causalité vs corrélation}

\begin{warningbox}[title=Limite importante]
Les valeurs SHAP indiquent des associations statistiques, \textbf{non des relations causales}.
Deux features fortement corrélées peuvent se "partager" l'importance SHAP sans que l'une
soit plus causale que l'autre. L'interprétation doit rester prudente.
\end{warningbox}

% ══════════════════════════════════════════════════════════════════════════════
\chapter{Dashboard et API}
% ══════════════════════════════════════════════════════════════════════════════

\section{Dashboard Streamlit}

Le dashboard constitue l'interface décisionnelle orientée \textbf{utilisateur métier}
(entraîneur, médecin du sport). Il est développé avec \textbf{Streamlit} et intègre des
graphiques interactifs \textbf{Plotly}.

\begin{table}[H]
\centering
\begin{tabularx}{\textwidth}{lX}
\toprule
\textbf{Page} & \textbf{Contenu} \\
\midrule
Vue d'ensemble & KPIs clés, répartition des classes, histogramme volume, matrice de corrélations 7 jours \\
Comparaison des modèles & Tableau comparatif stylisé, radar chart multi-dimensions, courbes PR/ROC, matrices de confusion, validation croisée \\
Prédiction & Formulaire 7 jours (onglets par jour), jauge de risque animée, bannière d'alerte colorée avec recommandations cliniques \\
Interprétabilité SHAP & Bar chart importance globale, beeswarm plot, classement top 20 features, guide de lecture \\
Simulation & Sliders profil coureur, courbe de risque en fonction de la réduction de charge (what-if) \\
\bottomrule
\end{tabularx}
\caption{Pages du dashboard RunSafe AI}
\end{table}

\section{API REST FastAPI}

\begin{table}[H]
\centering
\begin{tabular}{lll}
\toprule
\textbf{Méthode} & \textbf{Endpoint} & \textbf{Description} \\
\midrule
GET & \texttt{/health} & Vérification de l'état du service et des modèles \\
GET & \texttt{/model-info} & Modèles disponibles, seuils, noms d'affichage \\
POST & \texttt{/predict} & Prédiction du risque à partir de données 7 jours \\
\bottomrule
\end{tabular}
\caption{Endpoints de l'API FastAPI}
\end{table}

L'architecture respecte le principe \textbf{Front / API / Modèle} : le dashboard appelle
l'API pour les prédictions, indépendamment du chargement local des modèles.

Les schémas Pydantic (\texttt{api/schemas.py}) assurent la validation des entrées :
types, bornes, champs requis, messages d'erreur HTTP explicites.

% ══════════════════════════════════════════════════════════════════════════════
\chapter{Conclusion et perspectives}
% ══════════════════════════════════════════════════════════════════════════════

\section{Bilan}

Ce projet a permis de concevoir un système d'IA complet pour la prédiction de blessures
sportives, couvrant l'ensemble du pipeline professionnel : acquisition, nettoyage,
feature engineering, modélisation comparative, évaluation, interprétabilité, dashboard
et API.

Les principaux apprentissages sont :
\begin{itemize}
  \item \textbf{Le déséquilibre de classes} est le défi central : accuracy et ROC-AUC
        seuls sont trompeurs ; la PR-AUC est la métrique de choix.
  \item \textbf{Le Deep Learning n'est pas universellement supérieur} : le MLP est ici
        inférieur au Random Forest, faute de données suffisantes pour la classe minoritaire.
  \item \textbf{L'ACWR} (ratio charge aiguë/chronique) est la feature la plus importante,
        ce qui valide la pertinence métier du feature engineering.
  \item \textbf{L'interprétabilité SHAP} transforme un modèle opaque en outil de conseil
        actionnable pour le staff sportif.
\end{itemize}

\section{Perspectives}

\begin{itemize}
  \item \textbf{Données longitudinales} : prendre en compte l'historique complet de
        l'athlète (et non uniquement les 7 derniers jours) avec des modèles séquentiels
        (LSTM, Transformer).
  \item \textbf{Personnalisation} : entraîner des modèles individuels par athlète ou par
        profil (niveau, spécialité, morphologie).
  \item \textbf{Monitoring en production} : détecter la dérive des données (data drift)
        et ré-entraîner automatiquement le modèle.
  \item \textbf{Intégration IoT} : coupler avec des données GPS et cardiaques en temps
        réel pour des alertes plus précises.
\end{itemize}

% ══════════════════════════════════════════════════════════════════════════════
\chapter*{Références}
\addcontentsline{toc}{chapter}{Références}
% ══════════════════════════════════════════════════════════════════════════════

\begin{enumerate}
  \item Gabbett, T.J. (2016). The training-injury prevention paradox: should athletes
        be training smarter and harder? \textit{British Journal of Sports Medicine}, 50(5), 273--280.

  \item Lundberg, S.M., \& Lee, S.I. (2017). A unified approach to interpreting model
        predictions. \textit{Advances in Neural Information Processing Systems}, 30.

  \item Breiman, L. (2001). Random forests. \textit{Machine Learning}, 45(1), 5--32.

  \item Chen, T., \& Guestrin, C. (2016). XGBoost: A scalable tree boosting system.
        \textit{Proceedings of the 22nd ACM SIGKDD}, 785--794.

  \item Pedregosa, F., et al. (2011). Scikit-learn: Machine learning in Python.
        \textit{Journal of Machine Learning Research}, 12, 2825--2830.

  \item He, H., \& Garcia, E.A. (2009). Learning from imbalanced data.
        \textit{IEEE Transactions on Knowledge and Data Engineering}, 21(9), 1263--1284.

  \item Kaggle Dataset — Sports Injury Prediction (anonymisé, usage académique).
        \url{https://www.kaggle.com/datasets/vlad1986/sports-injury-prediction}
\end{enumerate}

% ══════════════════════════════════════════════════════════════════════════════
\appendix
\chapter{Structure du projet}
% ══════════════════════════════════════════════════════════════════════════════

\begin{lstlisting}[language=bash, caption=Arborescence du projet]
projet_2/
|-- data/raw/           # Dataset CSV brut
|-- src/
|   |-- preprocessing.py    # Pipeline preprocessing + feature engineering
|   |-- models.py           # 5 modeles ML/DL
|   +-- evaluation.py       # Metriques, SHAP, cross-validation
|-- dashboard/app.py        # Interface Streamlit (5 pages)
|-- api/
|   |-- main.py             # API FastAPI (/predict /health /model-info)
|   +-- schemas.py          # Schemas Pydantic
|-- notebooks/01_EDA_and_modeling.ipynb
|-- figures/                # Courbes, matrices de confusion, SHAP
|-- results/                # CSV comparatifs, seuils, features SHAP
|-- saved_models/           # Modeles serialises (.pkl)
|-- train.py                # Script d'entrainement principal
|-- run_cv.py               # Cross-validation 5-Fold
+-- requirements.txt
\end{lstlisting}

\chapter{Instructions de démarrage}

\begin{lstlisting}[language=bash, caption=Commandes de démarrage]
# 1 - Installation des dependances
pip install -r requirements.txt

# 2 - Entrainement (tous les modeles, ~5 min)
python train.py

# 3 - Cross-validation (~15 min, optionnel)
python run_cv.py

# 4 - Dashboard Streamlit
python -m streamlit run dashboard/app.py
# -> http://localhost:8501

# 5 - API REST (optionnel)
uvicorn api.main:app --reload --port 8000
# -> http://localhost:8000/docs
\end{lstlisting}

\end{document}
